\documentclass[aspectratio=169]{beamer}
\usepackage{amsmath}
\usepackage{booktabs}
\usepackage{threeparttable}
\definecolor{navyblue}{RGB}{26,43,76}
\definecolor{accentred}{RGB}{176,58,46}
\usetheme{default}
\setbeamertemplate{navigation symbols}{}
\setbeamertemplate{footline}[frame number]
\setbeamercolor{frametitle}{bg=navyblue,fg=white}
\setbeamerfont{frametitle}{series=\bfseries}
\setbeamercolor{itemize item}{fg=accentred}
\setbeamertemplate{itemize item}{\color{accentred}\textbullet}

\begin{document}

\begin{frame}{Research question}
  \vfill
  \begin{center}
    {\LARGE\bfseries
    Did offices that hired more people in response to drafting have more
    former staff return anywhere in the civil service after the war?}
  \end{center}
  \vfill
  \begin{itemize}
    \item Unit: 1944 destination office $\times$ team $\times$ position cell.
    \item Staffing response: internal workers plus external new hires.
    \item Postwar return: a 1944 draftee reappears in any available Tokyo
          civil-service directory during 1946--1955.
  \end{itemize}
  \vfill
\end{frame}

\begin{frame}{Empirical design}
  \scriptsize
  \textbf{Equations}
  \begin{align*}
    R_{op} &= \beta S_{op}+\gamma B_{op}+\alpha_s+\lambda_p+\varepsilon_{op}
      && \text{(OLS)},\\
    S_{op} &= \pi D_{op}+\gamma B_{op}+\alpha_s+\lambda_p+u_{op}
      && \text{(first stage)},\\
    R_{op} &= \beta \widehat S_{op}+\gamma B_{op}+\alpha_s+\lambda_p+\nu_{op}
      && \text{(second stage)}.
  \end{align*}

  \vspace{-0.7em}
  \textbf{Exact variable definitions}
  \begin{itemize}
    \item \textbf{Unit $(o,p)$:} a 1944 office ID $\times$ team (kakari)
          $\times$ normalized position-title cell.
    \item \textbf{$R_{op}$ --- returning workers:} number of men drafted from
          the cell in 1944 who appear at least once anywhere in the Tokyo civil
          service in an available 1946--1955 directory; each man counts once.
    \item \textbf{$S_{op}$ --- staffing response:} 1944 internal workers plus
          external new hires. Internal workers are retention in the same section,
          transfers from another section in the same department, and transfers
          from another department. A new hire is in his first observed directory
          year; opening-year office observations are excluded from this category.
          \textbf{This is a gross response: draftees are not subtracted.}
    \item \textbf{$D_{op}$ --- instrument:} number of male workers in the cell
          explicitly marked drafted in the 1944 directory.
    \item \textbf{$B_{op}$ --- baseline employment:}
          $\log(1+\text{number of distinct male workers previously observed in
          the same office ID $\times$ team $\times$ position})$.
    \item \textbf{Fixed effects:} 1944 section (bureau $\times$ section) and
          normalized position title. Standard errors are clustered by office ID.
  \end{itemize}
\end{frame}

\begin{frame}{Offices with larger staffing responses had more returnees}
  \centering
  \small
  \begin{threeparttable}
    \begin{tabular}{lccc}
      \toprule
      & OLS & IV & First stage \\
      \midrule
      Staffing response
        & 0.217$^{***}$ & 0.476$^{***}$ & \\
        & (0.038) & (0.067) & \\
      No. drafted workers
        & & & 0.848$^{***}$ \\
        & & & (0.087) \\
      \midrule
      Baseline employment control & Yes & Yes & Yes \\
      Section and position FE & Yes & Yes & Yes \\
      Observations & 654 & 654 & 654 \\
      Offices & 168 & 168 & 168 \\
      Cluster-robust first-stage $F$ & & 94.8 & \\
      \bottomrule
    \end{tabular}
    \begin{tablenotes}[flushleft]
      \tiny
      \item \textit{Notes:} The outcome is the number of 1944 draftees from the
      cell observed anywhere in the postwar Tokyo civil service. Staffing
      response is internal workers plus new hires; draftees are not subtracted.
      Standard errors, clustered by office, are in parentheses. $^{***}p<0.01$.
      \item \textit{Interpretation:} The IV estimate implies that one additional
      draft-induced staffing response predicts 0.48 additional postwar
      returnees. Because returnees are drawn from the drafted cohort, draft
      count may affect returns directly; the exclusion restriction is therefore
      a substantive limitation.
    \end{tablenotes}
  \end{threeparttable}
\end{frame}

\end{document}
